\section{Protocolo y Hoja de Registro de las Pruebas de Banco}
\apendiceindice{E}{Protocolo y Hoja de Registro de las Pruebas de Banco}

\subsection{Propósito}
Las pruebas de banco comprueban, antes del piloto, que el lector y el
servidor validan correctamente los códigos y que una lectura completa no
tarda más que la revisión manual. Se hacen con el prototipo armado fuera del
acceso, con estudiantes del equipo como usuarios.

\subsection{Pruebas de Lectura}
Se mide el tiempo de una lectura completa, desde que la pantalla con el
código queda frente a la cámara hasta que el servidor responde y el
controlador libera el torniquete. Se combinan tres condiciones de
iluminación (luz natural directa, sombra y luz artificial), dos niveles de
brillo de pantalla (50\,\% y 100\,\%) y dos distancias entre la pantalla y la
cámara (15 y 30~cm), lo que da 12 condiciones. En cada una se hacen 30
lecturas, para un total de 360. Cada lectura se registra como exitosa si
el código se reconoce al primer intento, en no más de cinco segundos.

\subsection{Pruebas de Validación}
Se comprueba que el servidor acepte solo los códigos válidos. Cada caso de la
lista se prueba al menos cinco veces:
\begin{enumerate}
	\item Código vigente de un estudiante dado de alta (debe aceptarse).
	\item Código vencido.
	\item Código ya utilizado (repetido).
	\item Código alterado en uno o más caracteres.
	\item Código con firma inválida.
	\item Código de un estudiante dado de baja.
	\item Desfase de reloj dentro de la tolerancia (debe aceptarse).
	\item Desfase de reloj fuera de la tolerancia.
\end{enumerate}

\subsection{Criterios de Aceptación}
El prototipo pasa al piloto si cumple tres condiciones: el 100\,\% de los
casos 2 a 6 y 8 se rechaza y los casos 1 y 7 se aceptan; al menos el
95\,\% de las lecturas de cada condición resulta exitosa (meta provisional,
que el equipo puede ajustar con los resultados); y el tiempo medio de una
lectura completa es menor que el tiempo medio de validación de la línea base
(Apéndice~C).

\subsection{Hoja de Registro}
\noindent Fecha: \linea[3cm] \quad Responsable: \linea[5cm]\\[0.5em]
Versión del firmware: \linea[3cm]

\medskip
\noindent
{\small
	\begin{tabular}{|p{2.4cm}|p{1.4cm}|p{1.6cm}|p{1.1cm}|p{2.2cm}|p{1.8cm}|p{2.1cm}|}
		\hline
		\textbf{Iluminación} & \textbf{Brillo} & \textbf{Distancia} & \textbf{Núm.} & \textbf{Tiempo (s)} & \textbf{¿Exitosa?} & \textbf{Notas} \\
		\hline
		                     &                 &                    & 1             &                     &                    &                        \\[0.6em] \hline
		                     &                 &                    & 2             &                     &                    &                        \\[0.6em] \hline
		                     &                 &                    & 3             &                     &                    &                        \\[0.6em] \hline
		                     &                 &                    & \dots         &                     &                    &                        \\[0.6em] \hline
		                     &                 &                    & 30            &                     &                    &                        \\[0.6em] \hline
	\end{tabular}}

\medskip
\noindent
{\small
	\begin{tabular}{|p{6cm}|p{2cm}|p{2.6cm}|p{3.6cm}|}
		\hline
		\textbf{Caso de validación} & \textbf{Intentos} & \textbf{Resultado esperado} & \textbf{Resultado obtenido} \\
		\hline
		1. Código vigente           &                   & Aceptado                    &                             \\ \hline
		2. Código vencido           &                   & Rechazado                   &                             \\ \hline
		3. Código repetido          &                   & Rechazado                   &                             \\ \hline
		4. Código alterado          &                   & Rechazado                   &                             \\ \hline
		5. Firma inválida           &                   & Rechazado                   &                             \\ \hline
		6. Estudiante dado de baja  &                   & Rechazado                   &                             \\ \hline
		7. Desfase dentro de la tolerancia & & Aceptado &                             \\ \hline
		8. Desfase fuera de la tolerancia  & & Rechazado &                            \\ \hline
	\end{tabular}}
